\begin{musica}{Chamateia}{Popular açoreana}{Tuna Académica da Universidade Lusída do Porto}
No ber\ac{D}{ço} que a ilha encerra
Bebo as \ac{G}{ri}mas deste \ac{D}{can}to,
No mar \ac{Bm}{al}to desta \ac{F#m}{te}rra
Nada a \ac{G}{ra}z\ac{A}{ão} do meu \ac{D}{pran}to.

Mas no terreiro da vida
O jantar serve de ceia
E mesmo a dor mais sentida
Dá lugar à chamateia.

\nota{Refrão:}
Oh meu \ac{Bm}{bem},
Oh chamar\ac{F#m}{ri}ta,
Meu a\ac{Am}{len}to, vai e \ac{Em}{vem},
Vou em\ac{Gm}{bar}car nesta \ac{D}{dan}ça
Sapa\ac{A7}{teia}, oh meu \ac{D}{bem}.

Se a sapateia não der
P’ra acalmar minh’alma inquieta,
Estou p’r’o que der e vier
Nas voltas da chamarrita.

Chamarrita, sapateia
Eu quero é contradizer,
O aperto desta bruma
Que às vezes me quer vencer.

\nota{Refrão}
\nota{D: xx0232}
\nota{Em: 022000}
\nota{F#m: 244222}
\nota{G: 320033}
\nota{A: x02220}
\nota{Am: x02210}
\nota{Bm: x24432}

\posicoes{D: xx0232, Em: 022000, F#m: 244222, G: 320033, A: x02220, Am: x02210,
  Bm: x24432}
\end{musica}
